\documentclass[conference]{IEEEtran}
\usepackage{amsmath,amssymb}
\usepackage{booktabs}
\usepackage{graphicx}
\usepackage{url}
\usepackage{hyperref}
\title{Error Propagation in Multi-Step Agentic AI Workflows for Research Automation: A Controlled Simulation Study}
\author{\IEEEauthorblockN{[AUTHOR NAME]}
\IEEEauthorblockA{[DEPARTMENT / INSTITUTION]\\ \texttt{[EMAIL OR STUDENT ID, IF REQUIRED]}}}
\begin{document}
\maketitle

\begin{abstract}
Multi-step research automation pipelines pass intermediate outputs between dependent stages, creating a mechanism by which an early error can affect later processing. This paper analyzes a four-stage simulation comparing independent errors with conditional downstream propagation. Recalculation of the saved multi-seed CSV (24,000 rows across 30 seeds, 100 trials per condition and base rate) gives increases in mean stage errors of 0.107, 0.207, and 0.331 at base error rates of 0.1, 0.2, and 0.3, respectively. The binary any-error rate is identical between conditions because workflow failure is defined as at least one stage error and the conditions share matched random seeds and trials. These results are an analysis of saved simulation outputs, not an independent reproduction of the data-generating run and not evidence about deployed language-model agents. The paper reports the simulation rule, data checks, paired comparisons, interpretation boundaries, and remaining reproducibility requirements.
\end{abstract}

\begin{IEEEkeywords}
agentic AI, error propagation, research automation, simulation, reproducibility, workflow reliability
\end{IEEEkeywords}

\section{Introduction}
Research automation increasingly involves a sequence of dependent activities, such as question formulation, literature retrieval, evidence extraction, synthesis, and report generation. When one stage produces an incorrect or incomplete intermediate artifact, later stages may reuse it. This motivates a focused question: under otherwise matched conditions, does a conditional downstream-error mechanism increase the number of errors accumulated across a multi-stage workflow relative to an independent-error model?

This study uses a four-stage simulation to isolate that mechanism. It does not evaluate a production agent, a real literature-review benchmark, or the factual correctness of generated research claims. The intended contribution is a small, inspectable controlled comparison that makes its assumptions and outcome definitions explicit. Any claim about real agentic systems requires a separate empirical validation using actual agent traces and independently assessed task outcomes.

\section{Research Question and Contributions}
The primary research question is: \emph{In a four-stage simulated research workflow, how does a conditional downstream error-propagation rule affect stage-error accumulation compared with a constant per-stage error probability?}

The planned contributions are:
\begin{itemize}
  \item a controlled comparison between independent errors and conditional downstream propagation;
  \item an evaluation procedure that validates raw records before producing summaries and plots; and
  \item an explicit analysis of why stage-error count and binary workflow-failure rate answer different questions.
\end{itemize}
These are contributions of experimental design and analysis, not a claim of publication-level novelty.

\section{Related Work}
Prior work motivates the study of intermediate-state reliability, but it does not directly establish the effect size measured by this simulation. Reasoning approaches include chain-of-thought prompting, self-consistency, and tree-structured search \cite{wei2022cot,wang2023selfconsistency,yao2023tot}. Agent frameworks extend this line of work by combining reasoning, action, memory, reflection, or multi-agent collaboration, as illustrated by ReAct, Reflexion, Toolformer, ToolLLM, AutoGen, and MetaGPT \cite{yao2023react,shinn2023reflexion,schick2023toolformer,qin2024toolllm,wu2024autogen,hong2024metagpt}. Generative Agents and Voyager further explore persistent memory, planning, and skill acquisition in interactive environments \cite{park2023generative,wang2023voyager}. These systems are relevant context, not direct replications of the controlled error-generation process studied here.

A second body of work concerns retrieval and evidence grounding. Retrieval-augmented generation combines model generation with retrieved passages \cite{lewis2020rag}; Self-RAG adds adaptive retrieval and self-reflection, while related work studies how retrieval and generation can be enabled or improved \cite{asai2024selfrag,gao2023enabling}. WebGPT explores browsing-assisted answer generation, and RePAcT examines interactive tool use for multi-step tasks \cite{nakano2021webgpt,dongre2025respact}. These approaches address grounding or tool-supported task completion, but they do not make the simulation's propagation parameters empirically valid.

Evaluation work also provides complementary perspectives. AgentBench, WebArena, Mind2Web, and GAIA assess agents across interactive or general-assistance tasks; SWE-bench evaluates software-engineering task resolution \cite{liu2024agentbench,zhou2024webarena,deng2023mind2web,mialon2024gaia,jimenez2024swebench}. FActScore evaluates factual precision at the level of atomic claims, while XAI-related work provides broader context for explainability and interpretability \cite{min2023factscore,dwivedi2023xai}. These resources measure different outcomes from the count of simulated stage errors. Accordingly, they motivate careful outcome definition and future empirical validation, but they do not validate this paper's programmed propagation rule or show that its numerical results generalize to deployed agentic research systems. Bibliographic metadata and claim-to-source alignment must still be checked against the primary sources before submission.

\section{Methodology}
\subsection{Workflow and Conditions}
The simulated workflow has four ordered stages. Two conditions are compared:
\begin{itemize}
  \item \textbf{Independent:} each stage draws its error according to the baseline per-stage error process without increasing later-stage error probability because of an earlier error.
  \item \textbf{Propagation:} an earlier error can conditionally increase the probability of errors at downstream stages.
\end{itemize}
The exact propagation function, random-number implementation, and trial-generation code must be cited from the actual source file or repository commit in the final version. If the code does not match these descriptions, revise this section to describe the implementation rather than the intention.

\subsection{Experimental Settings}
The saved multi-seed file contains 24,000 rows: 30 seeds (2026–2055), 100 trials per seed, condition, and base error rate, two conditions, and four base rates (0.0, 0.1, 0.2, 0.3). The generator in the Phase III notebook uses four ordered stages. In the independent condition, each stage uses the base error probability p. In the propagation condition, the probability for a stage becomes min(1, p + 0.30) when the immediately preceding stage had an error; otherwise it remains p. The same seed is reused for the two conditions, and rows can be matched by seed, trial, and base rate. This is a synthetic simulation, not an execution trace from a language model or a real research task.

\subsection{Outcomes}
The primary outcome is the mean number of stage errors per workflow. The secondary outcome is the fraction of workflows with at least one stage error. For the secondary outcome, a Wilson 95\% confidence interval may be reported if the underlying records are independent Bernoulli trials and the trial definition is appropriate. If paired seeds or trials are used, add a paired analysis for the primary comparison rather than treating all observations as independent.

The binary failure indicator is defined as
\[
F = \mathbb{1}(E > 0),
\]
where $E$ is the number of stage errors in a workflow. This metric records whether any error occurred, not how many later-stage errors followed the first error.

\section{Experimental Setup and Validation}
The evaluator should validate required columns, missing values, allowed ranges, consistency between the failure flag and stage-error count, duplicate trial keys, and condition/rate alignment before summarizing results. Any malformed or ambiguous rows must be reported, not silently removed. The analysis should preserve the raw input file and export a traceable summary table and figures.

\textbf{Data audit status:} the uploaded raw CSVs were inspected and recalculated. The multi-seed file has 24,000 rows, with 3,000 rows per condition and base rate; required fields are present, error rates are within [0,1], stage-error counts are nonnegative, failure flags equal the indicator of at least one stage error, and the (seed, trial, error-rate, model) keys are unique. The saved validation report also records a PASS. These checks validate data structure and summary arithmetic; they do not independently regenerate the CSVs. The included Python evaluator must still be executed in a documented environment and its output retained as part of the final reproduction record.

\section{Results}
Table~\ref{tab:results} summarizes the recalculated multi-seed records. Each condition has $n=3000$ observations at each base rate. Paired differences are calculated by matching rows on seed, trial, and base error rate. The displayed confidence intervals are approximate 95\% normal intervals for the mean paired difference (mean difference $\pm 1.96$ standard errors); they are descriptive and should be confirmed with the final analysis script.

\begin{table*}[t]
\centering
\caption{Recalculated stage-error counts and any-stage-error rates from the saved multi-seed CSV.}
\label{tab:results}
\small
\resizebox{\textwidth}{!}{%
\begin{tabular}{rrrrrl}
\toprule
Base rate & $n$/condition & Independent & Propagation & Difference (95\% CI) & Any-error rate (both) \\
\midrule
0.0 & 3000 & 0.0000 & 0.0000 & 0.0000 [0.0000, 0.0000] & 0.0000 \\
0.1 & 3000 & 0.3847 & 0.4917 & 0.1070 [0.0933, 0.1207] & 0.3360 \\
0.2 & 3000 & 0.7887 & 0.9953 & 0.2067 [0.1886, 0.2247] & 0.5787 \\
0.3 & 3000 & 1.1790 & 1.5100 & 0.3310 [0.3092, 0.3528] & 0.7447 \\
\bottomrule
\end{tabular}%
}
\end{table*}

The paired records show nonnegative differences for all matched trials; this is consistent with the shared-random-number construction and the propagation rule, and must be interpreted as a property of this simulation design. The binary any-error rate is identical across conditions because the outcome is defined as $F=\mathbb{1}(E>0)$: after a workflow has its first error, additional errors increase the count $E$ but cannot change $F$ from 1 to a larger value.

\section{Discussion}
If raw-data validation confirms that the propagation condition has a higher mean stage-error count, the result would be consistent with the simulation's intended mechanism: an early error can raise the chance of additional downstream errors. The size of this difference should be described using the verified estimates and uncertainty, not by visual inspection alone.

Equal any-error workflow-failure rates would not necessarily contradict increased stage-error accumulation. Under $F=\mathbb{1}(E>0)$, a workflow becomes a failure as soon as its first error occurs. Additional downstream errors can increase $E$ without changing $F$ from zero to one. Therefore, stage-error count is more sensitive to propagation severity, while the binary indicator measures the probability that a workflow experiences at least one error.

No conclusion should be drawn about real-world research automation reliability from this toy simulation alone. The simulation does not directly measure hallucination, citation validity, source relevance, synthesis quality, or end-to-end research correctness.

\section{Limitations and Threats to Validity}
First, this is a simulated process; its conclusions depend on the chosen error-generation rule. Second, a four-stage workflow may not represent the varied structures and feedback loops of deployed agents. Third, a stage error is a simplified event and may not capture severity, detectability, reversibility, or semantic dependence. Fourth, the saved CSV has been structurally audited and its summary statistics recalculated, but the full evaluator and data-generation workflow have not yet been independently rerun in a clean environment. Fifth, the paired confidence intervals rely on matching seeds and trials and should be reproduced by the final script. Finally, the related-work bibliography is drawn from the repository's existing bibliography and must be checked against primary sources and the citation audit before submission.

\section{Conclusion}
This paper specifies a controlled simulation for examining whether conditional downstream error propagation increases the number of errors accumulated in a multi-stage research workflow. The key methodological point is that accumulated error count and the probability of any workflow error are distinct outcomes. The numerical observations remain provisional until recomputed from raw trial-level data. The study should be presented as a mechanism-isolation simulation, with real-agent validation and verified literature synthesis left as necessary future work.

\section*{Data and Code Availability}
Repository: \url{https://github.com/gauravsu617/Awesome-agentic-ai-error-propagation}. The multi-seed input is tracked at \path{atr-submission/results/predictions/multi_seed_results.csv}. From the repository root, run the evaluator with that input and write outputs to \path{atr-submission/results/phase_iv_recomputed}; the exact command is documented in \path{docs/EXPERIMENT_DATA_AUDIT.md}. Preserve the evaluator output and generated plots. The source notebook uses a synthetic error-generation rule; this is not real-agent trace data. Confirm the intended author metadata and distribution rights before submission.

\section*{Acknowledgment of AI Assistance}
AI assistance was used to help structure and draft this manuscript and supporting documentation. The author must verify every source, inspect the code, rerun the analysis, confirm all numerical values, and disclose AI use according to institutional policy. This statement should be adjusted to match the actual tools and work performed.

\bibliographystyle{IEEEtran}
\bibliography{references}
\end{document}
